\begin{table}[t]
\centering
\caption{Robustness of hallucination-like ranking under alternative criteria. Values are percentages of all utterances; samples with trigram or four-gram repetition are excluded from the hallucination-like numerator.}
\label{tab:hallucination_robustness}
\small
\setlength{\tabcolsep}{5pt}
\begin{tabular}{lrrrrr}
\toprule
\textbf{Criterion} & \textbf{Base} & \textbf{RR} & \textbf{RU} & \textbf{UR} & \textbf{UU} \\
\midrule
Mean WER / mean Qwen3 & 17.00 & 16.50 & 16.88 & 15.95 & 16.91 \\
Top-25% WER / top-25% Qwen3 & 10.36 & 10.15 & 10.25 & 10.49 & 10.56 \\
WER > 0.5, Qwen3 > 0.6 & 1.97 & 1.98 & 2.18 & 5.03 & 2.09 \\
WER > 0.5, Qwen3 > 0.7 & 1.83 & 1.81 & 2.02 & 4.67 & 1.95 \\
\bottomrule
\end{tabular}
\vspace{0.5em}
\begin{minipage}{0.95\linewidth}
\footnotesize
\textit{Notes.} Hallucination-like outputs are treated as an operational diagnostic category: fluent or language-model-plausible hypotheses with low reference grounding. Repetition/oscillation is analyzed separately by excluding samples with elevated trigram or four-gram repetition from the hallucination-like numerator.
\end{minipage}
\end{table}
